\section{Rational Gamma reductions and explicit antiderivatives}\label{tsc:gamma}
At Stieltjes index zero, the antisymmetric first Stieltjes constant has a
finite Gamma reduction. We include a proof to make the all-depth evaluation
self-contained.

\begin{proposition}[Rational antisymmetric first Stieltjes value]\label{tsc:rational-D}
For integers $0<p<q$, with $a=p/q$,
\begin{align}\label{tsc:rational-D-eq}
 \gamma_1(1-a)-\gamma_1(a)
 ={}&\pi\cot(\pi a)\bigl(\gamma+\log(2\pi q)\bigr)\notag\\
 &-2\pi\sum_{r=1}^{q-1}\sin(2\pi pr/q)\log\Gamma(r/q).
\end{align}
Coprimality of $p,q$ is not required.
\end{proposition}
\begin{proof}
Write $z=e(a)$ and
$P(u)=\Gamma(u)(2\pi)^{-u}$. Subtract the two Hurwitz Fourier expansions:
\[
 \zeta(1-u,1-a)-\zeta(1-u,a)
      =2\ii P(u)\sin(\pi u/2)
                    \bigl(\Li_u(z)-\Li_u(z^{-1})\bigr).
\]
Now $2\ii P(u)\sin(\pi u/2)=\pi\ii(1-Au+O(u^2))$.
For real $u$, the difference of polylogarithms is $2\ii\Im\Li_u(z)$.
Comparing coefficients of $u$ gives
\[
 \gamma_1(1-a)-\gamma_1(a)
       =\pi A\cot(\pi a)-2\pi\Im\Li_0^{[1]}(z),
\]
since $\Im\Li_0(z)=\tfrac12\cot(\pi a)$.
Differentiate \eqref{tsc:finite-DFT} at zero. The constant part of Lerch's
formula vanishes in the finite sum, so
\[
 \Im\Li_0^{[1]}(z)=
     \sum_{r=1}^{q-1}\sin(2\pi pr/q)\log\Gamma(r/q)
                -\frac{\log q}{2}\cot(\pi a).
\]
Substitution proves \eqref{tsc:rational-D-eq}.
\end{proof}
This rational reflection formula is a consequence of classical Hurwitz
Fourier and Gamma identities. Its use here is to finish a whole family of
reflected product evaluations, rather than to claim a new isolated formula
for $\gamma_1(p/q)$.

\begin{theorem}[All-depth rational Gamma evaluation]\label{tsc:rational}
Let $a,b_1,\ldots,b_N$ be pairwise distinct rational circle points. Choose
$q$ with $d_j=\{b_j-a\}=p_j/q$, $0<p_j<q$. Then
\begin{align}\label{tsc:all-depth-gamma}
 &\FP\int_0^1 C_0(x-a)\prod_{j=1}^NK(x-b_j)\dd x\notag\\
 &\quad=\sum_{j=1}^N A_j(\mathbf b)
   \left\{K(d_j)\bigl(\gamma+\log(2\pi q)\bigr)
      -2\pi\sum_{r=1}^{q-1}\sin(2\pi p_jr/q)\log\Gamma(r/q)
      -\frac{\pi^2}{2}\right\}.
\end{align}
Thus this arbitrary-depth subclass has a finite expression in
$\gamma,\pi$, logarithms, and log-Gamma values at rational arguments,
with algebraic trigonometric coefficients.
\end{theorem}
\begin{proof}
Insert \eqref{tsc:rational-D-eq} into $I_0(d)$ in
\eqref{tsc:Iexamples}, then into Theorem~\ref{tsc:all-depth}.
At rational angles, the trigonometric coefficients belong to cyclotomic
number fields. No arithmetic independence statement is used.
\end{proof}
For even $N$, the final constants sum to zero by
\eqref{tsc:cot-residue-sum}. For odd $N$, their sum is
$-\pi^{N+1}\sin(N\pi/2)/2$. The parity effect checks the normalization.

\subsection{A cubic identity with only interior principal values}
\begin{theorem}[Quarter--half cubic digamma integral]\label{tsc:explicit-cubic}
The following principal value exists and satisfies
\begin{align}\label{tsc:explicit-cubic-eq}
 \boxed{\displaystyle
 \PV\int_0^1\psi(x)
     \cot\!\bigl(\pi(x-\tfrac14)\bigr)
     \cot\!\bigl(\pi(x-\tfrac12)\bigr)\dd x
  =\gamma+4\log2+3\log\pi-4\log\Gamma(\tfrac14).}
\end{align}
The principal values are symmetric at $1/4$ and $1/2$. The singularity at
zero is removable; the integrand tends to $-\pi$ there.
\end{theorem}
\begin{proof}
Here $A_1=K(-1/4)=-\pi$ and $A_2=K(1/4)=\pi$. The partial fraction
identity is
\[
 K(x-\tfrac14)K(x-\tfrac12)
        =-\pi^2-\pi K(x-\tfrac14)+\pi K(x-\tfrac12).
\]
Since $C_0=-\psi$, Theorem~\ref{tsc:all-depth} shows that the left side of
\eqref{tsc:explicit-cubic-eq} is
\[
 \frac{I_0(1/4)-I_0(1/2)}\pi
      =\frac{\gamma_1(3/4)-\gamma_1(1/4)}\pi.
\]
At $p/q=1/4$, Proposition~\ref{tsc:rational-D} gives
\[
 \frac{\gamma_1(3/4)-\gamma_1(1/4)}\pi
  =\gamma+\log(8\pi)-2\log\Gamma(1/4)+2\log\Gamma(3/4).
\]
Euler reflection, $\Gamma(1/4)\Gamma(3/4)=\pi\sqrt2$, reduces this to
the claimed right side.

Finally, $\psi(x)=-1/x+O(1)$,
$\cot(\pi(x-1/2))=-\pi x+O(x^3)$, and
$\cot(\pi(x-1/4))=-1+O(x)$. Their product tends to $-\pi$.
Consequently the one-sided finite-part notation needed in the general
theorem disappears in this specialization: only the two stated interior
principal values remain.
\end{proof}
The common value is
\[
 1.63190394589720479122328086909541565458453405168\ldots.
\]
The evaluation is proved by the preceding identities; the decimal is a
separate quadrature check.

\subsection{A log-Gamma cotangent transform}
\begin{proposition}[An ordinary primitive of the reflected kernel]\label{tsc:Gcot}
For $0<a<1$,
\begin{align}\label{tsc:Gcot-eq}
 \PV\int_0^1\log\Gamma(x)K(x-a)\dd x
   =\frac{\zeta^{[2]}(0,a)+\zeta^{[2]}(0,1-a)}2
        +\frac{\pi^2}{2}\left(a-\frac12\right).
\end{align}
The endpoint at zero is ordinarily integrable. In particular,
\begin{equation}\label{tsc:Gtan}
 \PV\int_0^1\log\Gamma(x)\tan(\pi x)\dd x
    =\frac{\log2\log(2\pi)+\tfrac12\log^22}{\pi}.
\end{equation}
\end{proposition}
\begin{proof}
Let the left side of \eqref{tsc:Gcot-eq} be $L(a)$.
By \eqref{tsc:Uone}, the periodic distributional derivative of
$\log\Gamma(x)-\ell/2$ is $-C_0$. Differentiation of the separated
pairing, followed by distributional integration by parts, gives
$L'(a)=-I_0(a)$.
Differentiating $\partial_x\zeta(s,x)=-s\zeta(s+1,x)$ at $s=0$, with
the spectral pole already canceled, gives
\[
 \partial_x\zeta^{[2]}(0,x)=2\gamma_1(x).
\]
Hence the derivative of the proposed right side is
$\gamma_1(a)-\gamma_1(1-a)+\pi^2/2=-I_0(a)$.
Both sides have mean zero in $a$: this is the zero mean of $K$ on the
left and of the Hurwitz jets at zero on the right. Their constant difference
is therefore zero.

At $a=1/2$, use $K(x-1/2)=-\pi\tan\pi x$ and
$\zeta(s,1/2)=(2^s-1)\zeta(s)$. Twice differentiating at zero yields
$\zeta^{[2]}(0,1/2)=-\log2\log(2\pi)-\tfrac12\log^22$, which proves
\eqref{tsc:Gtan}.
\end{proof}

More generally, an explicit pointwise antiderivative of every $I_m(a)$ in
\eqref{tsc:I-formula} is obtained by replacing $\gamma_n(a)$ with
\[
 Q_n(a)=\frac{(-1)^{n+1}}{n+1}\zeta^{[n+1]}(0,a),\qquad Q_n'(a)=\gamma_n(a),
\]
replacing $\gamma_n(1-a)$ with $-Q_n(1-a)$, and integrating the constant
term linearly. This produces a finite Hurwitz-jet antiderivative at every
Stieltjes index. A chosen base-point constant or a mean-zero periodic
normalization must then be specified separately.
